\chapter{Repo tour \& project constraints}

\section{What you're looking at}
This repository is intentionally small, but it isn't a single ``runtime''.
It is a collection of cooperating subsystems, each with its own contract and its own tests.

At a high level, you should think of Relax.js as three products living in one repo:

\begin{enumerate}
  \item \textbf{Relax VDOM runtime} (the teaching-friendly baseline) under \texttt{src/}.
  \item \textbf{HRBR runtime} (experimental, compiled fast-path) under \texttt{runtime/}.
  \item \textbf{Compiler} (experimental transform from TSX/JSX subset to HRBR blocks) under \texttt{compiler/}.
\end{enumerate}

\subsection*{Diagram: repository architecture at a glance}
\begin{center}
\begin{tikzpicture}[
  node distance=10mm and 12mm,
  box/.style={draw, rounded corners, align=left, inner sep=6pt, fill=RelaxLight},
  arrow/.style={-Latex, thick}
]
\node[box] (src) {\textbf{VDOM runtime}\\\texttt{src/*}\\mount/patch/destroy\\components, events};
\node[box, right=of src] (tests) {\textbf{Tests-as-spec}\\\texttt{src/\_\_tests\_\_/*}\\\texttt{runtime/\_\_tests\_\_/*}\\\texttt{compiler/\_\_tests\_\_/*}};
\node[box, below=of src] (runtime) {\textbf{HRBR runtime}\\\texttt{runtime/*}\\signals, scheduler\\blocks, hydration};
\node[box, below=of runtime] (compiler) {\textbf{Compiler}\\\texttt{compiler/*}\\JSX\,$\rightarrow$\,BlockDef};
\node[box, right=of runtime] (bench) {\textbf{Benchmarks + docs}\\\texttt{benchmarks/*}\\\texttt{docs/*}};

\draw[arrow] (src) -- node[midway, above] {verified by} (tests);
\draw[arrow] (runtime) -- (tests);
\draw[arrow] (compiler) -- (tests);

\draw[arrow] (compiler) -- node[midway, left] {emits defs consumed by} (runtime);
\draw[arrow] (src) -- node[midway, left] {can embed} (runtime);
\draw[arrow] (bench) -- (runtime);
\draw[arrow] (bench) -- (src);
\end{tikzpicture}
\end{center}

The key idea: tests are the stable external contract, while the implementation is split into a flexible VDOM baseline and an optional compiled runtime.

Everything else exists to support one of these (docs, benchmarks, build tooling).

\subsection{Directory map (first pass)}
This repository is split into a few major areas:

\begin{itemize}
  \item \texttt{src/}: the Relax.js Virtual DOM runtime.
  \item \texttt{runtime/}: an experimental runtime (HRBR) based on signals, scheduling, and compiled blocks.
  \item \texttt{compiler/}: a compiler/transform that converts a JSX subset into HRBR block definitions.
  \item \texttt{benchmarks/}: a browser harness to compare performance.
  \item \texttt{docs/}: design notes and whitepapers.
\end{itemize}

Two more folders matter when you're thinking about correctness:

\begin{itemize}
  \item \texttt{src/\_\_tests\_\_/}: the VDOM runtime spec.
  \item \texttt{runtime/\_\_tests\_\_/} and \texttt{compiler/\_\_tests\_\_/}: the HRBR + compiler spec.
\end{itemize}

\section{The constraint that shapes everything: DOM correctness in jsdom}
Almost every implementation choice is constrained by one hard rule:

\begin{quote}
The runtime must be correct in a headless DOM (jsdom), and tests are the contract.
\end{quote}

That affects:

\begin{itemize}
  \item Attribute vs property behavior (inputs/selects are tricky in jsdom).
  \item Event binding and cleanup.
  \item Diff algorithms (keyed moves must preserve node identity).
\end{itemize}

\subsection{Why jsdom is a meaningful constraint}
jsdom is \emph{not} a full browser, and that is exactly why it is useful here.
If your runtime depends on unspecified browser behavior or timing quirks, your test suite will become flaky.

Instead, tests in this repo focus on:

\begin{itemize}
  \item DOM structure (\texttt{innerHTML})
  \item identity (whether an element is reused or replaced)
  \item cleanup (event listeners and references are removed)
\end{itemize}

This shapes file boundaries: many helpers exist to make DOM updates deterministic and assertable.

\section{How to read code in this repo (recommended path)}
When you see a source file, immediately find its tests. That's the spec.

Recommended reading order:

\begin{enumerate}
  \item \texttt{src/index.ts} (public exports)
  \item \texttt{src/h.ts} (VNode model)
  \item \texttt{src/mount-dom.ts} (mount)
  \item \texttt{src/patch-dom.ts} (patch)
  \item \texttt{src/destroy-dom.ts} (cleanup)
  \item \texttt{src/component.ts} (components)
\end{enumerate}

\subsection{A second reading path: follow a user action}
The fastest way to build intuitions is to follow a single user interaction end-to-end.
For example: ``typing in an input and clicking a button adds a todo''.

You can trace that across:

\begin{itemize}
  \item event normalization: \texttt{src/events.ts}
  \item component event binding: \texttt{src/dispatcher.ts} and \texttt{src/component.ts}
  \item scheduling: \texttt{src/scheduler.ts} (\texttt{nextTick()})
  \item patching: \texttt{src/patch-dom.ts}
\end{itemize}

Doing this with the tests open beside you is the theme of this book.

\section{A concrete example: createApp() is tiny on purpose}
	exttt{createApp()} exists so the user has a simple lifecycle API.
The full file is short: \texttt{src/app.ts}.

Key behavior:

\begin{itemize}
  \item It creates a root VNode with \texttt{h(RootComponent, props)}.
  \item It mounts it once.
  \item It can unmount and fully cleanup.
\end{itemize}

\subsection{Code reference}
See \texttt{src/app.ts}. The central calls are \texttt{mountDOM(...)} and \texttt{destroyDOM(...)}.

\subsection{Test reference (what proves the contract)}
The high-level application lifecycle is specified by \texttt{src/\_\_tests\_\_/app.test.ts}.
This test file shows a full mount-to-interaction flow:

\begin{itemize}
  \item It mounts the app.
  \item It waits for async work with \texttt{await nextTick()}.
  \item It simulates user input by setting \texttt{input.value} and dispatching an \texttt{input} event.
  \item It simulates clicks via \texttt{button.click()}.
  \item It asserts DOM structure with \verb|singleHtmlLine`...`|.
\end{itemize}

That single test is a great ``reading harness'' for learning how Relax wires components + events + patching.

\section{Project constraint: small API surface}
Relax exports only a few primitives (h, fragments, components, scheduler).
Everything else is internal. That keeps the library small and also makes the tests more important.

\subsection{Where to see the supported surface area}
The public exports are in \texttt{src/index.ts} and mirrored by the \texttt{exports} map in \texttt{package.json}.

In other words:

\begin{itemize}
  \item if it's not in \texttt{src/index.ts}, it's not a stable public API
  \item tests are free to depend on internals, but users shouldn't
\end{itemize}

This book leans on that design: we can explain the internals without promising them as user-facing contracts.

\section{Tests: how this repo measures correctness}
This repository uses Vitest with a jsdom environment.
The configuration lives in \texttt{vitest.config.js}:

\begin{itemize}
  \item test patterns: \texttt{**/\_\_tests\_\_/**/*.test.ts(x)}
  \item reporters: verbose (good for learning)
  \item environment: jsdom
\end{itemize}

\subsection{The single most important helper: singleHtmlLine}
Many tests assert against HTML strings.
If you try to compare raw multi-line template literals to \texttt{innerHTML}, whitespace will fight you.

This repo solves that in \texttt{src/\_\_tests\_\_/utils.ts} with \texttt{singleHtmlLine()}:

\begin{itemize}
  \item it collapses whitespace
  \item it removes spaces between tags (for example, \texttt{>\textbackslash s+<} becomes \texttt{><})
  \item it trims
\end{itemize}

This is a micro-lesson in test design: you want assertions that are strict about structure, but forgiving about irrelevant formatting.

\subsection{Where to learn semantics fastest (test tour)}
If you're new to the codebase, these tests teach the contracts in the right order:

\begin{itemize}
  \item \texttt{src/\_\_tests\_\_/mount-dom.test.ts}:
    mount behavior, fragment behavior, saving \texttt{vdom.el} references.
  \item \texttt{src/\_\_tests\_\_/patch-dom.test.ts}:
    identity reuse vs replacement, attribute updates, keyed children, and edge cases.
  \item \texttt{src/\_\_tests\_\_/destroy-dom.test.ts}:
    cleanup and idempotency, especially ``remove event listeners''.
  \item \texttt{src/\_\_tests\_\_/component.test.ts}:
    component lifecycle, props/state updates, and event binding to the host component.
  \item \texttt{src/\_\_tests\_\_/jsx-react-syntax.test.tsx}:
    TSX compatibility contract (className/onClick/fragments/components).
\end{itemize}

Each of these will be revisited later in the book when we dive into the implementation files they protect.

\section{Mini-exercise}
Open these tests and write down what they prove (one sentence each). Then locate the implementation file they exercise.

\begin{itemize}
  \item \texttt{src/\_\_tests\_\_/app.test.ts}
  \item \texttt{src/\_\_tests\_\_/component.test.ts}
  \item \texttt{src/\_\_tests\_\_/mount-dom.test.ts}
  \item \texttt{src/\_\_tests\_\_/destroy-dom.test.ts}
\end{itemize}

If you can answer these, you're ready for Part II.

\section{Diagram: a "follow the tests" reading workflow}
When you don't know where to start, start from the tests and work backward to the implementation.

\begin{center}
\begin{tikzpicture}[
  node distance=9mm and 12mm,
  box/.style={draw, rounded corners, align=center, inner sep=6pt, fill=RelaxLight},
  arrow/.style={-Latex, thick}
]
  \node[box] (fail) {A behavior question\\("What should happen?")};
  \node[box, right=of fail] (test) {Find a test\\\texttt{src/\_\_tests\_\_/mount-dom.test.ts}};
  \node[box, right=of test] (assert) {Read assertions\\(DOM structure/identity)};
  \node[box, below=of assert] (impl) {Jump to implementation\\\texttt{src/*} / \texttt{runtime/*}};
  \node[box, left=of impl] (confirm) {Confirm contract\\(invariants, edge cases)};

  \draw[arrow] (fail) -- (test);
  \draw[arrow] (test) -- (assert);
  \draw[arrow] (assert) -- (impl);
  \draw[arrow] (impl) -- (confirm);
  \draw[arrow] (confirm) -- (test);
\end{tikzpicture}
\end{center}

This loop is also how you safely add features: you extend the spec (tests) and then evolve the implementation.

\section{Online resources}
This chapter is about reading \emph{this repo}, but a few external references help explain why these constraints and structures are common.

\begin{itemize}
  \item DOM living standard (event/bubbling/capturing model):
    \href{https://dom.spec.whatwg.org/}{WHATWG DOM Standard}.
  \item Event handling basics and differences between properties vs attributes:
    \href{https://developer.mozilla.org/en-US/docs/Web/API/EventTarget/addEventListener}{MDN: addEventListener},
    \href{https://developer.mozilla.org/en-US/docs/Web/API/HTMLElement}{MDN: HTMLElement}.
  \item Testing DOM in Node:
    \href{https://vitest.dev/guide/}{Vitest Guide},
    \href{https://github.com/jsdom/jsdom}{jsdom}.
  \item Module packaging and exports maps:
    \href{https://nodejs.org/api/packages.html#packages_exports}{Node.js docs: \texttt{exports}}.
\end{itemize}

For formal citations, you can map key links into the backmatter \texttt{References} list in \texttt{book/main.tex}.
